% ============================================================
% sec/02stage1.tex - Etapa 1 (21/08 - 17/09)
% ============================================================
\section{Etapa 1: estado del arte y dise\~no}
\label{sec:stage1}

\notebox[Objetivo de la etapa]{Elegir el tema, revisar la literatura en tres
ejes, auditar el corpus y entregar un art\'iculo IEEE preliminar con la
pregunta de investigaci\'on, las hip\'otesis y el protocolo experimental.
Entrega: 6 de septiembre de 2026. Peso: 30\,\%.}

\subsection{Resumen de ramas}

\begin{table}[H]
  \centering
  \caption{Ramas de la Etapa 1 y su resultado.}
  \label{tab:stage1summary}
  \footnotesize
  \begin{tabularx}{\textwidth}{@{}l P{2.0cm} P{2.6cm} Y@{}}
    \toprule
    \textbf{Rama} & \textbf{Fechas} & \textbf{Responsables} & \textbf{Resultado}\\
    \midrule
    \code{TrackConfirmation} & 21/08 & Adriel & Repositorio, licencia y reglas de Git\\
    \code{00TrackConfirmation} & 21/08--30/08 & Equipo & Track~A y tema aprobados; corpus auditado\\
    \code{01LiteratureReview} & 30/08--05/09 & Daniel (A), Adriel (B), Sebasti\'an (C) & Revisi\'on por ejes y bibliograf\'ia verificada\\
    \code{02PaperAndDelivery} & 05/09--17/09 & Equipo & Art\'iculo preliminar, declaraci\'on de IA y entrega\\
    \bottomrule
  \end{tabularx}
\end{table}

\subsection{Bit\'acora}

\jentry{2026-08-21}{Primera reuni\'on y creaci\'on del repositorio}%
       {Adriel, Daniel, Sebasti\'an}{TrackConfirmation}

El equipo se reuni\'o por primera vez por llamada (Fig.~\ref{fig:reunion1}).
Se revisaron conjuntos de datos p\'ublicos en Kaggle, entre ellos las
repeticiones del PTCG AI Battle Challenge publicadas por The Pok\'emon
Company. Ese mismo d\'ia Adriel cre\'o el repositorio con su licencia, su
\code{.gitignore} y las reglas de ramas y commits.

\begin{figure}[H]
  \centering
  \safeimage[width=0.85\textwidth]{../21-8-26_PrimeraReunionDelEquipo.png}
  \caption{Primera reuni\'on del equipo (21/08/2026), revisando conjuntos de
  datos en Kaggle.}
  \label{fig:reunion1}
\end{figure}

\decision{Organizar el repositorio en tres niveles de ramas
(\code{master}, \code{Stage<N>} y ramas de trabajo numeradas) y usar
Conventional Commits.}{La autor\'ia individual se califica por el historial de
commits, y una rama por bloque de trabajo hace visible qui\'en hizo qu\'e y
cu\'ando.}{Ramas personales y una rama \code{develop}: mezclaban trabajos y
ocultaban la etapa a la que pertenec\'ia cada cambio.}

\milestone{Repositorio operativo.}{PR \#1 y commit ``docs(repo): add
branching and commit conventions''.}

% ------------------------------------------------------------

\jentry{2026-08-26}{Confirmaci\'on del tema}%
       {Adriel, Sebasti\'an}{00TrackConfirmation}

En una sesi\'on de revisi\'on (Fig.~\ref{fig:confirmacion}) se cerr\'o el
tema: estimar la probabilidad de victoria en el Pok\'emon TCG desde el estado
de la partida, como clasificaci\'on tabular.

\begin{figure}[H]
  \centering
  \safeimage[width=0.40\textwidth]{../26-8-26_ConfirmacionDelTema.png}
  \caption{Sesi\'on de revisi\'on en la que se confirm\'o el tema
  (26/08/2026).}
  \label{fig:confirmacion}
\end{figure}

\decision{Track~A (aprendizaje cl\'asico sobre datos tabulares) con el corpus
del PTCG AI Battle Challenge.}{El corpus es p\'ublico, grande y tiene
etiqueta clara (qui\'en gana). El problema de evaluaci\'on que plantea, miles
de estados correlacionados por partida, es justo el que el Track~A pide
tratar con rigor: fuga, baselines y protocolo.}{Una primera direcci\'on de
aprendizaje por refuerzo (planificaci\'on de rutas de cobertura en terreno
desconocido). Se abandon\'o tras la reevaluaci\'on del equipo y la
retroalimentaci\'on del profesor.}

\aiuse{Claude}{explorar direcciones del proyecto y probar si la pregunta era
falsable}{la elecci\'on del tema, del \textit{track} y la decisi\'on de dejar
la direcci\'on de refuerzo fueron del equipo}

% ------------------------------------------------------------

\jentry{2026-08-30}{Auditor\'ia del corpus y base de investigaci\'on}%
       {Adriel, Daniel, Sebasti\'an}{00TrackConfirmation}

En una sesi\'on de trabajo de los tres (Fig.~\ref{fig:etapa1}) se cerr\'o la
confirmaci\'on del \textit{track} y se fij\'o la base de investigaci\'on. La
auditor\'ia del corpus analiz\'o diez d\'ias entre el 17 de junio y el 22 de
julio: 52\,085 partidas y 7\,038\,378 observaciones, unos 180~GiB, sin errores
de lectura. Tambi\'en midi\'o que, en cada paso, la vista del jugador que
espera es 98.8\,\% id\'entica a su vista anterior.

\begin{figure}[H]
  \centering
  \safeimage[width=0.40\textwidth]{../30-8-26_DesarrolloDeLaEtapa1.png}
  \caption{Sesi\'on de trabajo de la Etapa 1 (30/08/2026).}
  \label{fig:etapa1}
\end{figure}

\decision{Agrupar por partida toda partici\'on, validaci\'on cruzada e
intervalo.}{La auditor\'ia midi\'o que una partici\'on aleatoria por fila
colocar\'ia pr\'acticamente todas las partidas a ambos lados de la divisi\'on:
el modelo se evaluar\'ia sobre estados casi id\'enticos a los que vio. La
decisi\'on se apoya en una medici\'on sobre el propio corpus, no en una
analog\'ia.}{Partici\'on aleatoria por fila. Solo se conserva como
demostraci\'on del sesgo, nunca para reportar resultados.}

\decision{Tomar solo la vista del jugador que debe actuar y describir cada
estado desde su perspectiva (propio, oponente y diferencia).}{La vista del
jugador que espera repite casi siempre la anterior y duplicar\'ia filas sin
informaci\'on nueva. La perspectiva propio/oponente evita que el modelo
aprenda ``la posici\'on 0 gana m\'as''.}{Usar ambas vistas por paso, o
columnas por posici\'on fija en la mesa.}

\decision{Declarar una lista de campos prohibidos: mano del oponente,
contenido de los premios, resultado, recompensas, nombres de agentes e
historial de jugadas.}{Son informaci\'on que el jugador no ve en ese instante
o que revela el desenlace. Usarlos ser\'ia fuga.}{Confiar en que el modelo
``no los usar\'ia'': la fuga se previene por construcci\'on, no por
inspecci\'on posterior.}

\milestone{Tema confirmado y corpus auditado.}{PR \#2 y commit ``docs(repo):
close track confirmation and establish the Stage 1 research base''.}

% ------------------------------------------------------------

\jentry{2026-08-30 a 2026-09-05}{Revisi\'on de literatura por ejes}%
       {Daniel (eje A), Adriel (eje B), Sebasti\'an (eje C)}{01LiteratureReview}

Cada integrante revis\'o un eje: el problema y su dominio (A), las t\'ecnicas
del \textit{track} (B) y el conjunto de datos (C). Se seleccionaron 21
art\'iculos arbitrados. Cada uno resolvi\'o los DOI de sus fuentes y ley\'o
el texto completo antes de agregarlas a la bibliograf\'ia.

\decision{Una entrada solo entra al \code{.bib} principal si su DOI resuelve
y su contenido corresponde a lo que se le atribuye.}{Las herramientas de IA
proponen metadatos plausibles pero a veces falsos; una cita incorrecta
invalida el argumento que sostiene.}{Aceptar los metadatos que sugiere una
herramienta. Cuando un dato no pudo confirmarse, se dej\'o vac\'io en lugar
de rellenarlo con una suposici\'on.}

\decision{Usar como baseline de literatura el AUC de 0.78 a 0.80 del
desaf\'io AAIA'17 sobre Hearthstone.}{Es el antecedente m\'as cercano: un
juego de cartas por turnos con predicci\'on de victoria desde el estado.}{
Usar cifras de juegos de estrategia en tiempo real, cuyo estado y din\'amica
son muy distintos.}

\blocker{Varios candidatos eran \textit{preprints} sin arbitraje.}{Retras\'o
el cierre de la bibliograf\'ia.}{Se marcaron como tales y se excluyeron, o
se consult\'o con el profesor antes de incluirlos.}

\aiuse{Consensus, Claude y ChatGPT}{buscar candidatos, clasificarlos y
detectar \textit{preprints}}{cada DOI se resolvi\'o a mano y cada art\'iculo
se ley\'o}

\milestone{Revisi\'on de literatura integrada al borrador.}{PR \#3.}

% ------------------------------------------------------------

\jentry{2026-09-05 a 2026-09-06}{Ensamblaje del art\'iculo y revisi\'on final}%
       {Adriel, Daniel, Sebasti\'an}{02PaperAndDelivery}

Se integraron la introducci\'on, los trabajos relacionados, la tabla
comparativa, el baseline de literatura, la metodolog\'ia, el an\'alisis
t\'ecnico y las consideraciones \'eticas. La revisi\'on final dur\'o m\'as
de hora y media con los tres integrantes (Fig.~\ref{fig:revision}).

\begin{figure}[H]
  \centering
  \safeimage[width=0.30\textwidth]{../06-9-26_RevisionFinalEtapa1.png}
  \caption{Revisi\'on final del art\'iculo de la Etapa 1 (06/09/2026).}
  \label{fig:revision}
\end{figure}

\decision{Formular tres hip\'otesis refutables (H1 a H3) sin umbrales
num\'ericos fijos.}{Sin resultados previos sobre este corpus, cualquier
umbral ser\'ia arbitrario. Cada hip\'otesis declara qu\'e resultado la
refutar\'ia.}{Hip\'otesis del tipo ``el modelo funciona bien'', que no
pueden fallar.}

\decision{Una escalera de cinco modelos, B0 a B4, en la que cada nivel
responde una pregunta.}{B0 (probabilidad a priori) fija el piso; B1
(heur\'istica) pregunta si hace falta aprendizaje; B2 (regresi\'on
log\'istica) si la relaci\'on es lineal; B3 (bosque) si hay interacciones;
B4 (potenciaci\'on del gradiente) cu\'anto a\~nade la complejidad.}{Comparar
solo contra un modelo trivial, que no dice si la complejidad vale la pena.}

\decision{Tres escenarios de evaluaci\'on: partidas nuevas (A), d\'ias
posteriores (B) y mazos reservados (C).}{Cada uno responde una pregunta de
generalizaci\'on distinta, y se reportan por separado.}{Un solo conjunto de
prueba aleatorio, que no mide el cambio de la poblaci\'on ni la
memorizaci\'on de mazos.}

\milestone{Art\'iculo preliminar entregado.}{PR \#4, PR \#5 (\code{Stage1}
$\rightarrow$ \code{master}) y etiqueta \code{v1.0.0}.}

% ------------------------------------------------------------

\jentry{2026-09-06 a 2026-09-17}{Correcciones posteriores a la entrega}%
       {Adriel}{02PaperAndDelivery}

Se corrigieron figuras, la integraci\'on de las referencias y el t\'itulo de
los autores.

\milestone{Correcciones integradas.}{PR \#6 (etiqueta \code{v1.0.1}) y
PR \#7.}

\subsection{Decisiones consolidadas}

\begin{table}[H]
  \centering
  \caption{Decisiones t\'ecnicas de la Etapa 1.}
  \label{tab:stage1decisions}
  \footnotesize
  \begin{tabularx}{\textwidth}{@{}l P{3.8cm} Y P{2.2cm}@{}}
    \toprule
    \textbf{Fecha} & \textbf{Decisi\'on} & \textbf{Por qu\'e} & \textbf{Secci\'on del art\'iculo}\\
    \midrule
    26/08 & Track~A con el corpus del PTCG AI Battle Challenge & Datos p\'ublicos con etiqueta clara y un reto de evaluaci\'on real & Introducci\'on\\
    30/08 & Agrupar todo por partida & Medido: la partici\'on por fila mezcla casi todas las partidas & Metodolog\'ia, an\'alisis t\'ecnico\\
    30/08 & Una fila por estado del jugador que act\'ua, en perspectiva propia & Evita duplicados y el sesgo de posici\'on & Metodolog\'ia\\
    30/08 & Lista de campos prohibidos & Previene la fuga por construcci\'on & An\'alisis t\'ecnico\\
    05/09 & Baseline de literatura: AUC 0.78--0.80 (Hearthstone) & Antecedente m\'as cercano & Baseline de literatura\\
    06/09 & Hip\'otesis H1 a H3 refutables & Sin umbrales arbitrarios & Introducci\'on\\
    06/09 & Escalera B0--B4 y escenarios A, B, C & Cada nivel y escenario responde una pregunta & Metodolog\'ia\\
    \bottomrule
  \end{tabularx}
\end{table}

\subsection{Retrospectiva}

\begin{table}[H]
  \centering
  \caption{Retrospectiva de la Etapa 1.}
  \label{tab:stage1retro}
  \footnotesize
  \begin{tabularx}{\textwidth}{@{}P{4.0cm} Y@{}}
    \toprule
    \textbf{Pregunta} & \textbf{Respuesta}\\
    \midrule
    ¿Qu\'e funcion\'o? & Repartir la literatura por ejes y medir el corpus antes de dise\~nar el protocolo. La decisi\'on de agrupar por partida qued\'o respaldada por una medici\'on propia.\\
    ¿Qu\'e nos bloque\'o? & Verificar los metadatos de cada referencia contra la fuente primaria tom\'o m\'as tiempo del previsto.\\
    ¿Qu\'e estimamos mal? & El tiempo de ensamblaje del art\'iculo: la mayor parte del trabajo se concentr\'o en los \'ultimos dos d\'ias.\\
    ¿Qu\'e cambia para la Etapa 2? & Traducir la documentaci\'on al espa\~nol, reemplazar el plan \'agil por un plan de desarrollo por ramas y abrir una rama propia para los datos antes de entrenar.\\
    \bottomrule
  \end{tabularx}
\end{table}